% Part II -- Closed-loop surface facility conceptual design.
% All CO2 properties: Span-Wagner via CoolProp; well inventories from
% the HyperCO2 hydrostatic+geothermal initializer (n = 200 cells).

\part*{Part II --- Conceptual Design of the Closed-Loop Surface
Facility}
\addcontentsline{toc}{part}{Part II --- Closed-loop surface facility
conceptual design}

\section{II.1 Scope, method, and the two numbers that drive the design}

Part II is a preliminary engineering feasibility study for the surface
facility that will charge and discharge the 2,000-ft research well
with CO$_2$ --- repeatedly, with reproducible initial conditions, and
as a closed loop wherever that is defensible. It deliberately does not
assume the initial concept (high-pressure tank + choke + compressor)
is right; the configurations of \S II.4 are compared on computed
thermodynamics, and the recommendation in \S II.12 differs from the
initial concept in two essential ways. All properties are Span--Wagner
\citep{span1996} via CoolProp \citep{bell2014}; well inventories come
from the Part-I model's hydrostatic--geothermal initializer. Cost
statements are order-of-magnitude judgments, flagged as such; nothing
here substitutes for formal engineering review (\S II.10).

Two computed numbers dominate every configuration decision:
\begin{enumerate}
\item \textbf{The working inventory is $\sim$1.6--2.1 tonnes.} The
charged well (wellhead \SI{5.0}{MPa}, hydrostatic column, summer
geothermal profile) holds \SI{1742}{kg}; a residual
\SI{179}{kg} remains at a \SI{0.7}{MPa} wellhead, so one full
blowdown transfers $\approx$\SI{1563}{kg}. A cold, uniform
\SI{300}{K} column holds \SI{2064}{kg} (the Part-I stage-1 state), so
storage must be sized for the cold-column worst case,
$\approx$\SI{2.1}{t}, not the geothermal nominal.
\item \textbf{Gas-phase collection at low back pressure is
volumetrically impractical; liquid collection is trivial.} Holding the
transferred \SI{1563}{kg} as gas at \SI{300}{K} requires
\SI{122}{m^3} at \SI{0.7}{MPa}, \SI{84}{m^3} at \SI{1.0}{MPa}, and
still \SI{40}{m^3} at \SI{2.0}{MPa}; as refrigerated liquid
(\SI{2.0}{MPa}, \SI{-20}{\celsius}, $\rho = \SI{1032}{kg/m^3}$) it
occupies \SI{1.5}{m^3}. Any closed loop therefore ends in a
\emph{condensing} recovery train, and the only question is how fast
it must run (\S II.3, \S II.8).
\end{enumerate}

\section{II.2 CO$_2$ state space of the facility}

The facility operates between four thermodynamic anchors
(\cref{tab:anchors}): refrigerated liquid storage on the saturation
line near \SI{2}{MPa}/\SI{-20}{\celsius} (the industrial
``mini-bulk'' standard, chosen in \S II.5); the charged well
(\SI{5.0}{MPa} wellhead, \SI{5.75}{MPa} bottomhole, gas throughout by
the Part-I margin analysis); the low-pressure recovery side
(\SIrange{0.55}{1.0}{MPa} gas); and ambient-temperature liquid
(saturation \SI{6.43}{MPa} at \SI{25}{\celsius}), which the design
\emph{avoids} holding in any closed volume because of
thermal-expansion overpressure (\S II.10). Two phase-diagram
boundaries act as hard constraints: the critical point
(\SI{7.377}{MPa}, \SI{304.13}{K}) --- the summer surface temperature
of \SI{305.4}{K} sits \emph{above} $T_c$, so ``liquid at ambient''
does not exist on hot days and near-critical storage states have
violently temperature-sensitive densities (\S II.4, Configuration 3b)
--- and the triple point (\SI{0.518}{MPa}, \SI{216.59}{K}): any
depressurization path that crosses \SI{0.518}{MPa} while cold can
form solid CO$_2$, which is why the closed-loop operating mode holds
its collection pressure \emph{above} the triple pressure
(\S II.3, Scenario B) and why vent and relief lines, which do cross
it, are designed for dry ice (\S II.10; cf.\ the Part-I triple-point
discussion and \citealp{martynov2018}).

\begin{table}[htb]
\centering
\caption{Thermodynamic anchor states (Span--Wagner). The well column
is the Part-I hydrostatic--geothermal profile; storage and recovery
states are design choices justified in \S II.4--\S II.5.}
\label{tab:anchors}
\begin{tabular}{lcccl}
\toprule
State & $p$ [MPa] & $T$ [\si{\celsius}] & $\rho$ [\si{kg/m^3}] &
Phase \\
\midrule
Liquid storage (mini-bulk) & 2.0 & $-20$ & 1032 & sat.\ liquid \\
Storage vapor head & 2.0 & $-20$ & 52 & sat.\ vapor \\
Charged well, wellhead & 5.0 & 32.3 & 121 & gas \\
Charged well, bottomhole & 5.75 & 47.8 & 129 & gas \\
Blowndown well, wellhead & 0.7 & var. & 13 & gas \\
LP buffer (closed mode) & 0.55--1.0 & amb. & 10--19 & gas \\
Ambient liquid (avoided) & 6.43 & 25 & 710 & sat.\ liquid \\
\bottomrule
\end{tabular}
\end{table}

On a $p$--$h$ diagram the operating cycle is a loop: pumping raises
the stored liquid from \SI{2.0}{} to \SI{6.5}{MPa} almost vertically
(\SI{6.7}{kJ/kg}); the vaporizer carries it horizontally across the
dome and into superheat (\SI{285}{kJ/kg}); the charging valve drops
it isenthalpically into the well; blowdown expands the well inventory
down the (approximately isentropic, wall-heated) paths computed in
Part I; and the recovery compressor plus condenser close the loop
back to the liquid anchor (\SI{101}{kJ/kg} of work, \SI{449}{kJ/kg}
of heat rejection at \SI{-20}{\celsius}). Phase changes by process
step: storage --- two-phase by design (saturated vessel); charging
valve --- single-phase gas provided the vaporizer superheats
sufficiently (\S II.3-A); wellbore --- single-phase by the Part-I
margin analysis for the geothermal protocol, two-phase for cold
protocols (the Part-I model extension exists to quantify exactly
this); choke discharge --- flashing/two-phase possible, already
treated by the Part-I HEM choke; LP collection --- gas (kept above
the dew line at ambient temperature by pressure choice); compression
--- gas, with \SI{143}{\celsius} discharge temperature at ratio 3.1
(aftercooled); condensation --- two-phase by design.

\section{II.3 Operating scenarios}

\paragraph{Scenario A --- charging the well.} The well starts at
\SIrange{0.7}{1.0}{MPa} (post-experiment) and must reach
\SI{5.0}{MPa} wellhead with a prescribed temperature state. The
recommended sequence: cryogenic pump draws saturated liquid from
storage, raises it to \SIrange{6.0}{6.5}{MPa} (subcooled by the
pressure rise itself, eliminating cavitation concern at the pump
\emph{discharge}; inlet NPSH is provided by vessel liquid head and a
flooded-suction installation); an ambient-air or trim-electric
vaporizer takes it to \SIrange{305}{315}{K} superheated gas
(\SI{285}{kJ/kg}; \SI{142}{kW} at \SI{0.5}{kg/s}); a flow-control
valve admits it to the wellhead. Isenthalpic expansion from
\SI{6.5}{MPa}/\SI{310}{K} to \SI{5.0}{MPa} lands at \SI{21.5}{\celsius}
--- \SI{7.2}{K} above the \SI{5}{MPa} dew point of
\SI{14.3}{\celsius} --- so the superheat setpoint directly controls
the margin against condensation in the injection line; below
$\approx$\SI{303}{K} vaporizer outlet the margin vanishes. Charging
\SI{1563}{kg} at \SI{0.5}{kg/s} takes $\approx$\SI{52}{min};
``staged pressurization'' is simply this sequence run under
mass-flow control with soak periods, and needs no extra equipment.
The alternatives are inferior on computed grounds: \emph{direct
pressure-driven transfer from a gas vessel} requires either a
\SI{28}{m^3} vessel at \SI{6.5}{MPa} (isothermal \SI{305}{K}, to
deliver \SI{1563}{kg} without falling below \SI{5.2}{MPa}) or a
near-critical dense vessel whose deliverable mass collapses from
515 to \SI{149}{kg/m^3} of density swing between \SI{305}{} and
\SI{313}{K} tank temperature --- an un-reproducible, weather-sensitive
initial condition; \emph{direct compressor charging} from the
\SI{2}{MPa} storage vapor head needs a \SIrange{30}{50}{kW}-class
machine at ratio $>3$ with aftercooling, i.e.\ more capital and more
failure modes than a \SI{7}{kW} pump, to do the same job more slowly;
it also boils the storage liquid by vapor withdrawal, dragging the
tank cold.

\paragraph{Scenario B --- blowdown.} Two modes, selected by what the
experiment needs. \emph{Vented mode} (reference experiments,
model-validation campaigns at fixed atmospheric back pressure):
discharge through the 12.7 mm choke to a vertical vent stack;
reproduces the Part-I boundary condition exactly; loses the
transferred inventory ($\approx$\SI{1.6}{t}/run --- a consumable
cost, not a technical problem, but incompatible with
``routine closed-loop'' aspirations). \emph{Closed mode}: discharge
into the LP buffer held at \SIrange{0.55}{1.0}{MPa} by a
back-pressure control valve, with the recovery train condensing into
storage concurrently. Two physical facts make this mode respectable
rather than a compromise: while the choke is \emph{choked},
downstream pressure does not influence the wellbore transient at all,
and with $p_b = \SI{0.7}{MPa}$ the flow stays choked until the
wellhead falls to $\approx p_b/0.522 \approx \SI{1.34}{MPa}$
($\gamma \approx 1.43$ at the relevant state) --- i.e.\ the majority
of the transient is identical to the vented experiment, and only the
subcritical tail differs (quantifiable in advance with the Part-I
model, which takes $p_b$ as an input); and holding
$p_b > \SI{0.518}{MPa}$ keeps the entire downstream system above the
triple pressure, excluding dry-ice formation everywhere except the
relief/vent paths. The buffer cannot, at reasonable size, absorb a
full-rate 12.7 mm blowdown ($\dot m_0 \approx \SI{1.7}{kg/s}$,
\S II.1 item 2); closed mode therefore pairs naturally with
\emph{reduced chokes} (4--6 mm, $\dot m \lesssim \SI{0.3}{kg/s}$),
for which the recovery train keeps up in near--real time ---
precisely the orifice sizes of the \citet{hammer2022} experiments.
Full-rate closed operation would need the buffer sized by
$V = m_{\mathrm{transfer}}/\rho(p_{b,\max})$: \SI{84}{m^3} even
allowing the back pressure to end at \SI{1.0}{MPa}.

\paragraph{Scenario C --- recovery and reset.} The recovery
compressor takes buffer gas (\SI{0.7}{MPa}, ambient) to
\SI{2.2}{MPa} (single stage, ratio 3.1, \SI{101}{kJ/kg},
\SI{143}{\celsius} discharge), an aftercooler/condenser rejects
\SI{449}{kJ/kg} at the \SI{-20}{\celsius} storage level, and
condensate returns to the storage vessel. Turnaround is set entirely
by this train: \SI{1563}{kg} at \SI{0.1}{kg/s} is \SI{4.3}{h}
(\SI{10}{kW} compressor, \SI{45}{kW} refrigeration); at
\SI{0.2}{kg/s}, \SI{2.2}{h} (\SI{20}{kW}, \SI{90}{kW}). Thermal reset
of the \emph{well} is separate and slower: after a blowdown the
completion and near-wellbore formation are cold (Part-I stored-energy
results), and the geothermal initial condition recovers on the
formation's conductive time scale (hours; measurable with the DTS and
itself a conjugate-model validation target). Reproducible initial
conditions therefore favor a standardized soak protocol ---
e.g.\ charge, then hold overnight --- rather than faster surface
equipment.

\paragraph{Scenario D --- repeated cycles.} Repetition changes the
optimum: it is what disqualifies the single-gas-tank concept (state
drift, \S II.4) and what justifies the liquid anchor --- a saturated
vessel returns to the \emph{same} ($p$,$T$) after every cycle so long
as its refrigeration holds the setpoint, making storage state the one
part of the facility that is reproducible by construction. It also
elevates inventory accounting (load cells on the storage vessel,
\S II.9) from convenience to requirement: closed-loop leak-down is
otherwise invisible.

\section{II.4 Configurations compared}

\Cref{fig:pfd-configs} sketches Configurations 1--4;
\cref{fig:pfd-rec} details the recommended hybrid.
\emph{Configuration 1 --- single gas tank.} One vessel both sources
and receives. Fails quantitatively on both duties at once: as source
it needs \SI{28}{m^3} at \SI{6.5}{MPa} (above), as receiver its
pressure must sit \emph{below} the blowndown wellhead, after which
the same vessel cannot recharge the well without a compressor rated
for the full \SI{6.5}{MPa} discharge. State drift between cycles
(each experiment changes tank $p$, $T$, inventory) makes reproducible
initial conditions a control problem rather than a property of the
design. Rejected.
\emph{Configuration 2 --- HP gas tank + LP gas receiver +
compressor.} The initial concept. Feasible, but the HP tank carries
the Configuration-1 sizing burden and the LP receiver carries the
\SIrange{40}{120}{m^3} gas-collection burden; the inter-tank
compressor must then lift \SI{1.6}{t} across a ratio of
$\sim$9 per cycle. Every vessel is large, every transfer is
mechanical, and nothing in the loop is anchored to a reproducible
state. Dominated by Configuration 5 on all criteria except
conceptual familiarity.
\emph{Configuration 3 --- direct compressor charging, no HP tank.}
Smallest vessel count, but the charging compressor becomes the
facility: ratio $>3$ from storage-vapor suction, \SIrange{30}{50}{kW},
aftercooled, running for the entire charge, with the storage vessel
refrigeration fighting the boil-off its suction induces. Charging
rate, and hence the initial-condition protocol, is hostage to one
machine. The near-critical ``dense storage'' variant (3b) is rejected
for the \SI{305}{}--\SI{313}{K} density cliff computed in \S II.3-A.
\emph{Configuration 4 --- liquid storage + pump.} The charging side
of the recommended design: \SI{2.5}{m^3}-class refrigerated vessel,
\SI{7}{kW} pump, vaporizer. By itself it is open-loop (no recovery);
it answers the charging question but not the collection question.
\emph{Configuration 5 --- hybrid (recommended).} Liquid storage +
pump + vaporizer for charging; LP buffer + back-pressure valve for
collection; compressor + condenser for recovery; vent stack for
full-rate reference blowdowns. Each component exists because a
computed number demands it, and none duplicates another's function.
The ``more equipment'' caution in the brief is answered by what is
\emph{absent}: no HP gas storage, no second large receiver, no
charging compressor.

\begin{figure}[htb]
\centering
\begin{tikzpicture}[
  scale=0.82, every node/.style={transform shape},
  vessel/.style={draw, rounded corners=2pt, minimum width=1.5cm,
                 minimum height=0.85cm, align=center, font=\scriptsize},
  mach/.style={draw, circle, minimum size=0.72cm, align=center,
               font=\scriptsize, inner sep=1pt},
  well/.style={draw, minimum width=0.5cm, minimum height=1.5cm,
               align=center, font=\scriptsize},
  lbl/.style={font=\scriptsize\itshape},
  flow/.style={-{Stealth[length=2mm]}, thick}]
% Config 1
\begin{scope}
\node[vessel] (t1) at (0,0) {gas\\tank};
\node[well] (w1) at (3.1,0) {well};
\draw[flow] (t1) -- node[above,font=\tiny]{charge} (w1);
\draw[flow, dashed] (w1.south) |- ++(-1.5,-0.55) -| (t1.south)
  node[pos=0.25, below, font=\tiny]{blowdown + compressor};
\node[lbl] at (1.55,-1.9) {Config.\ 1: single gas tank};
\end{scope}
% Config 2
\begin{scope}[xshift=6.4cm]
\node[vessel] (hp) at (0,0.65) {HP\\tank};
\node[vessel] (lp) at (0,-0.95) {LP\\recv.};
\node[well] (w2) at (3.1,0) {well};
\node[mach] (c2) at (-1.5,-0.15) {C};
\draw[flow] (hp.east) -- node[above,font=\tiny]{charge} (w2.120);
\draw[flow] (w2.240) -- node[below,font=\tiny]{choke} (lp.east);
\draw[flow] (lp.west) -| (c2.south);
\draw[flow] (c2.north) |- (hp.west);
\node[lbl] at (1.2,-2.3) {Config.\ 2: two gas tanks};
\end{scope}
% Config 3
\begin{scope}[yshift=-4.6cm]
\node[vessel] (t3) at (0,0) {storage};
\node[mach] (c3) at (1.7,0) {C};
\node[well] (w3) at (3.4,0) {well};
\draw[flow] (t3) -- (c3);
\draw[flow] (c3) -- node[above,font=\tiny]{charge} (w3);
\draw[flow, dashed] (w3.south) |- ++(-2.0,-0.6) -| (t3.south)
  node[pos=0.25, below, font=\tiny]{blowdown};
\node[lbl] at (1.7,-1.95) {Config.\ 3: compressor charging};
\end{scope}
% Config 4
\begin{scope}[xshift=6.4cm, yshift=-4.6cm]
\node[vessel] (t4) at (0,0) {liquid\\storage};
\node[mach] (p4) at (1.45,0) {P};
\node[vessel] (v4) at (2.75,0) {vap.};
\node[well] (w4) at (4.35,0) {well};
\draw[flow] (t4) -- (p4); \draw[flow] (p4) -- (v4);
\draw[flow] (v4) -- (w4);
\draw[flow, dashed] (w4.south) -- ++(0,-0.55)
  node[below, font=\tiny]{vent (open loop)};
\node[lbl] at (2.1,-1.95) {Config.\ 4: liquid + pump (open)};
\end{scope}
\end{tikzpicture}
\caption{Block process-flow sketches of Configurations 1--4
(C = compressor, P = pump; dashed = discharge path). Configuration 5
is \cref{fig:pfd-rec}.}
\label{fig:pfd-configs}
\end{figure}

\begin{figure}[htb]
\centering
\begin{tikzpicture}[
  vessel/.style={draw, rounded corners=2pt, minimum width=1.9cm,
                 minimum height=1.0cm, align=center, font=\scriptsize},
  mach/.style={draw, circle, minimum size=0.8cm, align=center,
               font=\scriptsize, inner sep=1pt},
  valve/.style={draw, regular polygon, regular polygon sides=3,
                shape border rotate=270, minimum size=0.32cm,
                inner sep=0pt, font=\tiny},
  well/.style={draw, minimum width=0.6cm, minimum height=2.0cm,
               align=center, font=\scriptsize},
  lbl/.style={font=\scriptsize},
  flow/.style={-{Stealth[length=2.2mm]}, thick}]
\node[vessel] (stor) at (0,0)
  {liquid storage\\ \SI{2.0}{MPa}, $-20\,^{\circ}$C\\ 2.5--\SI{3}{m^3}, load cells};
\node[mach] (pump) at (2.6,0.75) {P};
\node[lbl, above=0.05 of pump] {\SI{7}{kW}};
\node[vessel] (vap) at (4.7,0.75) {vaporizer\\ 140 kW$_{\mathrm{th}}$};
\node[valve] (fcv) at (6.7,0.75) {};
\node[lbl, above=0.02 of fcv] {FCV};
\node[well] (well) at (8.6,0) {well\\ \SI{5}{MPa}};
\node[valve] (chk) at (8.6,-1.75) {};
\node[lbl, right=0.05 of chk] {choke 4--12.7 mm};
\node[valve] (bpv) at (6.6,-2.5) {};
\node[lbl, above=0.02 of bpv] {BPV \SI{0.7}{MPa}};
\node[vessel] (buf) at (4.4,-2.5)
  {LP buffer\\ 10--\SI{15}{m^3}, $\le$\SI{1}{MPa}};
\node[mach] (comp) at (2.2,-2.5) {C};
\node[lbl, above=0.05 of comp] {10--\SI{20}{kW}};
\node[vessel] (cond) at (0,-2.5)
  {condenser\\ 45--90 kW$_{\mathrm{th}}$ at $-20\,^{\circ}$C};
\node[lbl] (vent) at (10.6,-2.5) {vent stack};
\draw[flow] (stor.east) ++(0,0.3) -| (pump.south);
\draw[flow] (pump) -- (vap);
\draw[flow] (vap) -- (fcv);
\draw[flow] (fcv) -- (well.120);
\draw[flow] (well.south) -- (chk.north);
\draw[flow] (chk.south) |- (bpv.east);
\draw[flow] (bpv.west) -- (buf.east);
\draw[flow] (buf.west) -- (comp.east);
\draw[flow] (comp.west) -- (cond.east);
\draw[flow] (cond.west) -| ++(-0.5,1.2) |- (stor.west);
\draw[flow, dashed] (chk.south) ++(0,-0.2) -| ++(1.4,0) -- (vent.west)
  node[pos=0.5, above, font=\tiny]{vented mode};
\end{tikzpicture}
\caption{Recommended configuration (5): liquid-anchored closed loop
with a vented reference mode. FCV = flow-control valve, BPV =
back-pressure valve, P = cryogenic pump, C = recovery compressor.
Relief devices, isolation valves, and instrumentation (\S II.9--II.10)
omitted for clarity.}
\label{fig:pfd-rec}
\end{figure}

\section{II.5 Storage: how much, in what phase}

Sizing from the cold-column worst case with margin:
$m_{\mathrm{stor}} = 1.5 \times \SI{2.1}{t} \approx \SI{3.2}{t}$,
i.e.\ a \SI{3}{m^3}-class refrigerated vessel at
\SI{2.0}{MPa}/\SI{-20}{\celsius} --- squarely the commercial CO$_2$
``mini-bulk'' class (3--6 t, factory-insulated, packaged
refrigeration), available for purchase or rental from industrial-gas
suppliers, which is itself a strong argument for this anchor state:
the long-lead, code-stamped component becomes a catalog item. Storage
capacity should be based on \emph{total} inventory (not per-cycle
transfer): the closed loop holds one inventory that migrates between
well and storage, and the 1.5$\times$ factor covers heel, vapor head,
and commissioning losses. The storage pressure band
\SIrange{1.7}{2.3}{MPa} (\SI{-25}{} to \SI{-15}{\celsius}) trades
refrigeration duty against pump suction margin; \SI{2.0}{MPa} is the
conventional center. The minimum information to finalize sizing:
confirmed bore volume per completion string actually used
(\SI{13.9}{m^3} for the 7-5/8\,in string modeled in Part I), the
coldest planned initial condition, and the heel the pump can reach.

\section{II.6 Compressor versus pump: the division of labor}

The intuition that ``a compressor displaces CO$_2$ from one vessel to
another'' conflates three distinct operations with three different
costs. (i) \emph{Pumping liquid}: work
$w \approx v\,\Delta p/\eta$; with
$v = \SI{0.97e-3}{m^3/kg}$, \SIrange{2.0}{6.5}{MPa} and
$\eta = 0.65$, $w = \SI{6.7}{kJ/kg}$ --- \SI{7}{kW} at \SI{1}{kg/s},
discharge temperature essentially unchanged (\SI{-16.5}{\celsius}).
(ii) \emph{Compressing gas}: work scales with
$\int v\,dp$ at gas-phase specific volume, two orders of magnitude
larger; \SIrange{0.7}{2.2}{MPa} costs \SI{101}{kJ/kg}
($\eta_s = 0.70$) and lands at \SI{143}{\celsius}, demanding
aftercooling; ratios above $\sim$4 per stage force multistage
machines with intercooling. (iii) \emph{Pressure-driven transfer}:
free, but consumes the source vessel's state --- the \SI{28}{m^3}
sizing result of \S II.3-A is the price of using a gas vessel's own
pressure as the driver. The division of labor follows: the
\emph{pump} does the high-pressure lift (charging), because the fluid
is liquid there and the lift is cheap; the \emph{compressor} does the
low-pressure gathering (recovery), because the fluid is unavoidably
gas there and the ratio is modest; \emph{passive transfer} is used
only where it is free of state consequences (blowdown itself, and
gravity-fed pump suction). Both machines are justified --- but the
compressor is a \SI{10}{}--\SI{20}{kW} recovery unit, not the
\SI{50}{kW}-class charging machine of the initial concept. Per cycle
(\SI{1563}{kg}): pump 2.9 kWh$_{e}$; vaporizer
124 kWh$_{\mathrm{th}}$ (ambient-air units make most of this free);
recovery compression 44 kWh$_{e}$; condensation
195 kWh$_{\mathrm{th}}$ rejected at \SI{-20}{\celsius}
($\approx$100--130 kWh$_{e}$ at packaged-unit COP 1.5--2).
Liquid pumping beats gas compression for the charging duty by a
factor of $\sim$15 in specific work --- but the honest accounting
shows the cycle's real energy sits in the recovery refrigeration, not
in either machine.

\section{II.7 Equipment sizing methodology}

Governing relations, stated so the numbers above are reproducible and
re-sizable: vessel inventories from
$m = \rho(p,T)\,V$ with Span--Wagner $\rho$; the well inventory from
$m_{\mathrm{well}} = \int_0^L \rho\bigl(p(z),T(z)\bigr)A\,dz$ with
the Part-I hydrostatic--geothermal profile (evaluated numerically by
the model's initializer); pump power $\dot W = \dot m\,\Delta h_s/\eta$
with $\Delta h_s$ along the liquid isentrope; compressor power the
same along the gas isentrope, staged when $p_2/p_1 \gtrsim 4$;
vaporizer and condenser duties as enthalpy differences between the
anchor states; receiver volumes from $V = \Delta m/\rho(p_{\max},T)$;
charging and recovery times from $\Delta m/\dot m$; blowdown duration
from the Part-I model itself (e.g.\ wellhead-pressure half-life
\SI{494}{s} for the 12.7 mm choke at $C_d = 0.84$). Sensitivities
worth carrying into detailed design: all gas-side densities scale
$\sim p/T$, so receiver sizing is robust; the near-critical states
are the opposite (\S II.3-A) and should simply be avoided; vaporizer
duty varies $\pm$15\,\% across the storage band; compressor work is
insensitive to suction temperature but linear in $\ln(p_2/p_1)$.

\section{II.8 Energy and turnaround comparison}

Per full cycle, Configuration 5: $\approx$47 kWh$_{e}$ of machine
work plus 100--130 kWh$_{e}$ of refrigeration --- order
\$20--40 of electricity --- against $\approx$\SI{1.6}{t} of CO$_2$
(order \$300--600 at merchant prices, \emph{[commodity prices
indicative only]}) saved per run relative to venting. The loop pays
for its operating cost in avoided gas after a handful of runs; the
capital trade is the condenser/refrigeration train against the
otherwise-necessary receiver volume, and \S II.1 item 2 settles that
in the condenser's favor by an order of magnitude in vessel size.
Turnaround: charge \SI{\sim 1}{h}, experiment
\SIrange{0.3}{1.2}{h}, recovery \SIrange{2.2}{4.3}{h}
(compressor-limited), thermal soak of the formation to the geothermal
initial state: hours--overnight (conduction-limited, not
equipment-limited). One experiment per day with an overnight soak is
the natural cadence; a second daily run is possible at the cost of a
non-equilibrium (but model-characterizable) initial thermal state.

\section{II.9 Instrumentation and control}

Actively \emph{controlled}: charging mass flow (FCV on Coriolis
feedback), vaporizer outlet superheat ($\ge$\SI{10}{K} above the
injection-line dew point), closed-mode back pressure (BPV setpoint
\SIrange{0.55}{1.0}{MPa}), storage temperature (packaged
refrigeration), compressor suction pressure (recycle/unload).
\emph{Measured}: wellhead and choke-upstream/downstream $p$,$T$ (fast
transducers at the wellhead --- the Part-I wave analysis implies
$\gtrsim$\SI{100}{Hz} to resolve the expansion-wave arrival and
reflections); Coriolis mass flow on the charge line; blowdown mass
flow \emph{by inference} --- choke-model + upstream state, closed by
storage-vessel load cells and buffer $p$,$T$ inventory accounting ---
because a Coriolis meter in the discharge line would see flashing
two-phase flow, where its reading is not trustworthy; the existing
completion DTS for distributed temperature (the Part-I/Part-II
saturation-thermometer opportunity); tank levels by differential
pressure plus load cells; O$_2$/CO$_2$ area monitors. The control
philosophy mirrors the model architecture deliberately: every
controlled boundary of the facility (charge flow, back pressure,
storage state) is an input the Part-I model already accepts, which is
what makes the digital-twin coupling of \S II.11 a bookkeeping
exercise rather than a re-formulation.

\section{II.10 Safety (conceptual; formal review required)}

The hazards are the standard dense-CO$_2$ set, concentrated as
follows. \emph{Trapped-liquid overpressure}: any liquid-full segment
between two valves (pump discharge line, condensate return) generates
extreme pressure on warming; every such segment needs a thermal
relief valve --- this is the single most common CO$_2$-facility
design error. \emph{Two-phase/flashing relief}: relief devices on
liquid or dense service must be sized by HEM/omega-method two-phase
methods (API 520 Annex C; ISO 4126-10), not gas formulas --- the same
physics as the Part-I choke analysis, and \citet{hammer2022} shows
HEM sizing errs conservative (low) on flow, which for relief sizing
is the \emph{non}-conservative direction; a margin policy is
required. \emph{Dry ice}: vent and relief paths cross the triple
pressure by design; they must be short, straight, downward-sloping,
single-walled where blockage matters, and rated for
\SI{-78}{\celsius} impingement (low-temperature-toughness material
selection per B31.3); the closed-mode BPV floor of \SI{0.55}{MPa} is
itself a dry-ice control. \emph{Asphyxiation}: a full inventory
release is \SI{\sim 1.6}{t} of CO$_2$; outdoor, elevated vent
discharge, area monitors with alarm-at-\SI{0.5}{\%}/evacuate
levels, and attention to low-lying accumulation (CO$_2$ is
1.5$\times$ air density). \emph{Machine protection}: compressor
suction knockout (liquid carryover from the buffer on cold days),
pump minimum-flow recycle, high-discharge-temperature trips.
\emph{Codes}: ASME BPVC VIII vessels, B31.3 process piping, API
520/521 relief, with CGA/EIGA CO$_2$-specific guidance
\emph{[standard designations to be confirmed in formal review]}.
Pressure ratings: the charging side sees \SI{6.5}{MPa} + relief
margin; the well and wellhead govern themselves under existing
facility certification.

\section{II.11 Coupling to the wellbore model: the digital twin}

The surface equipment is lumped-parameter by construction --- every
vessel is a well-mixed mass/energy accumulator,
\begin{equation}
\frac{dM_i}{dt} = \sum_j \dot m_{ij},
\qquad
\frac{d(M_i u_i)}{dt} = \sum_j \dot m_{ij} h_{ij}
 + \dot Q_i + \dot W_i,
\label{eq:lumped}
\end{equation}
closed by Span--Wagner $(\rho, u) \to (p, T, x)$ --- the \emph{same}
UV-flash primitive recovery as the Part-I two-phase wellbore closure,
so one thermodynamic kernel serves the entire twin. Valves and chokes
use the Part-I HEM orifice flux (already implemented); the
charging FCV is the identical model run in reverse service; the
compressor is a map $\dot m(p_s, p_d, N)$ with isentropic efficiency;
the pump is $\dot m$-prescribed with \cref{eq:lumped} bookkeeping;
heat exchangers are effectiveness--NTU duties. Connecting piping
(tens of meters) has negligible volume next to the well
(\SI{13.9}{m^3}) and buffer; it needs distributed treatment only if
surface-line acoustics themselves become a measurement target. The
coupling to the distributed wellbore model is through exactly one
interface --- the wellhead boundary condition --- and the Part-I
choke closure already takes the downstream pressure as a parameter:
replacing the constant $p_b$ with the buffer state
$p_b = p_{\mathrm{buf}}(t)$ from \cref{eq:lumped}, integrated
explicitly at the wellbore time step (the buffer's fill time of
minutes against acoustic steps of milliseconds makes the explicit
coupling error negligible, by the same time-scale argument as the
Part-I conjugate thermal coupling), turns the existing solver into
the discharge half of the twin with no structural change. Mass and
energy conservation across the facility then reduces to summing
\cref{eq:lumped} over vessels plus the wellbore budgets the Part-I
code already closes to machine precision --- giving the twin a
built-in, auditable conservation check, and giving leak detection a
quantitative baseline.

\section{II.12 Recommendation}

\textbf{A. One tank or two?} Two vessels, but not the two of the
initial concept: one refrigerated \emph{liquid} storage
(3 t class) as the reproducible anchor, plus one modest LP gas
\emph{buffer} (10--\SI{15}{m^3}) as shock absorber --- not a second
full-inventory gas tank.
\textbf{B. Compressor for charging?} No. Charging is pump +
vaporizer; passive transfer alone also fails (\S II.3-A). The
compressor belongs on recovery, at one-fifth the power the charging
duty would demand of it.
\textbf{C. Liquid pump advantages?} Decisive: \SI{6.7}{} vs
\SI{101}{kJ/kg}, \SI{7}{kW} vs tens of kW, no discharge-temperature
management, catalog cryogenic hardware. Its obligations ---
flooded suction, NPSH, thermal-relief of liquid-full segments --- are
standard practice.
\textbf{D. Storage phase?} Saturated liquid at
\SI{2.0}{MPa}/\SI{-20}{\celsius}. Gas storage fails on volume,
ambient-temperature liquid on the critical point and
thermal-expansion hazard, near-critical dense fluid on
reproducibility.
\textbf{E. LP receiver rather than direct return?} Yes ---
structurally necessary, since the blowndown stream can neither enter
the \SI{2}{MPa} storage unaided nor be condensed at its instantaneous
peak rate; the buffer decouples the fast experiment from the slow
recovery train.
\textbf{F. Reset procedure?} Compressor--condenser transfer at
\SIrange{0.1}{0.2}{kg/s} (\SIrange{2.2}{4.3}{h}), concurrent with an
overnight formation soak for geothermal initial conditions.
\textbf{G/H. Best configuration, and easiest to model?} The same one
--- Configuration 5 --- and not coincidentally: its control variables
are exactly the model's boundary conditions (\S II.9, \S II.11).
\Cref{tab:matrix} summarizes.

\begin{table}[htb]
\centering
\small
\caption{Configuration comparison. ++ favorable, $-$ unfavorable,
$\circ$ neutral/conditional. Capital and operating costs are
order-of-magnitude judgments at concept level.}
\label{tab:matrix}
\begin{tabular}{lccccc}
\toprule
Criterion & C1 single & C2 two-gas & C3 compr. & C4 liq.\ open &
C5 hybrid \\
\midrule
Capital cost & $\circ$ & $-$ & $\circ$ & ++ & $\circ$ \\
Operating cost/energy & $-$ & $-$ & $-$ & $\circ$ & + \\
Equipment complexity & + & $-$ & $\circ$ & ++ & $\circ$ \\
Operational flexibility & $-$ & $\circ$ & $-$ & $\circ$ & ++ \\
Repeatability of ICs & $-$ & $-$ & $\circ$ & ++ & ++ \\
CO$_2$ retention (closed) & $\circ$ & + & $\circ$ & $-$ & ++ \\
Vessel volumes & $-$ & $-$ & + & ++ & + \\
Safety profile & $\circ$ & $\circ$ & $\circ$ & + & + \\
Ease of control & $-$ & $-$ & $-$ & + & + \\
Digital-twin suitability & $\circ$ & $\circ$ & $-$ & + & ++ \\
Future research headroom & $-$ & $\circ$ & $-$ & $\circ$ & ++ \\
\bottomrule
\end{tabular}
\end{table}

\paragraph{Development strategy.} \emph{Phase 1 (essential, enables
first experiments)}: liquid storage vessel, pump, vaporizer, charging
FCV + Coriolis, vent stack, wellhead fast instrumentation, relief and
gas-detection systems --- this is Configuration 4 and runs vented
reference blowdowns that validate the Part-I/Part-II models.
\emph{Phase 2 (closes the loop)}: LP buffer, BPV, recovery
compressor, condenser, load-cell inventory accounting --- converts
the same facility to Configuration 5 without rework, enabling
repeated closed-cycle campaigns and the reduced-choke experiments.
\emph{Optional/later}: second completion-string tie-ins, a
surface test section for choke-model experiments
(\citealp{hammer2022} geometry), slipstream impurity dosing.

\paragraph{Unresolved before detailed engineering.} Confirmed
completion string(s) and wellhead rating for repeated thermal
cycling; coldest intended initial condition (sets storage margin and
materials); site power available for the refrigeration train;
acceptable per-run venting policy (regulatory and institutional);
relief-study basis (full-bore charging-line failure vs choke stuck
open); merchant CO$_2$ supply logistics for first fill and makeup;
and the choke-hardware plan (fixed beans vs adjustable), which fixes
whether $C_d$ calibration (Part I, \S 10) happens on the well or on a
surface test section.

\section{II.13 References and precedents for the facility concept}

The closed-loop architecture has direct precedents in the
depressurization rigs reviewed in Part I: the SINTEF ECCSEL
depressurization facility \citep{munkejord2020, munkejord2021}
(instrumentation density and data practices worth copying: 100 kHz
pressure arrays, triplet thermocouples, open raw data), the
\citet{hammer2022} orifice test section (the natural template for a
surface choke-calibration leg), and the large-scale vented programs
(COOLTRANS; \citealp{cosham2012}) whose venting practice ---
vertical stacks, dry-ice-tolerant design --- informs the vented mode.
No published facility performs repeated \emph{wellbore-scale}
closed-loop CO$_2$ charge/discharge cycles; that is the same
experimental gap identified in Part I (\S 14H), and this facility is,
to the knowledge assembled here, the instrument that would fill it.
